% ---------------------------------------------------------------------------
%  The Four Quadrants of Quantum Machine Learning
%  Build:  latexmk -pdf main.tex      (or pdflatex, bibtex, pdflatex, pdflatex)
% ---------------------------------------------------------------------------
\documentclass[11pt,a4paper]{article}

\usepackage{qmlproject}
\qmlheader{The four quadrants of QML}

\title{\textbf{The Four Quadrants of Quantum Machine Learning}\\[0.4em]
  \large Concepts, methods, and where the QML Lab project fits}
\author{QML Lab project notes}
\date{\today}

\begin{document}
\maketitle

\begin{abstract}
\noindent
Quantum machine learning (QML) is not one field but four. A common taxonomy sorts work by two
questions: is the data (or the system being studied) classical or quantum, and is the computer that
learns from it classical or quantum? The answers give four quadrants, \CC, \CQ, \QC\ and \QQ. This
document explains each quadrant and the main concepts inside it, with the key equations and
references, and then maps the modules and research projects of the QML Lab onto the quadrants. Most
of the project lives in \CQ, the quadrant served by PennyLane and Qiskit; the open-ended projects
reach into \QC\ and \QQ, where the strongest current results and the clearest arguments for quantum
advantage are found.
\end{abstract}

\tableofcontents
\bigskip

% ===========================================================================
\section{Why a taxonomy?}
% ===========================================================================

``Quantum machine learning'' is used for very different activities: running a neural network on a
quantum computer, using a neural network to read out a quantum experiment, borrowing quantum
mathematics for classical algorithms, or learning directly from quantum states. These activities need
different hardware, rest on different assumptions and make different claims about speed-ups. Lumping
them together causes most of the confusion in the field.

The taxonomy used here goes back to A\"imeur, Brassard and Gambs \citeyearpar{aimeur2006}, who distinguished
learning settings by whether the learner and the learning task are classical or quantum. It was used
in the now-familiar $2\times2$ form by Dunjko, Taylor and Briegel \citeyearpar{dunjko2016,dunjko2018} and
popularised in the textbook of Schuld and Petruccione \citeyearpar{schuld2018,schuld2021book}. Broad reviews
of the field are refs. \citep{biamonte2017,dunjko2018,cerezo2021vqa}.

\begin{inshort}
Each label has two letters. The \emph{first} letter says what kind of data or system is being
studied; the \emph{second} says what kind of computer processes it. \CQ\ therefore means
\emph{classical data, quantum processor}.
\end{inshort}

% ===========================================================================
\section{Preliminaries: quantum states, measurement and ``quantum data''}
\label{sec:prelim}
% ===========================================================================

\paragraph{Qubits and states.}
A qubit is a unit vector in $\mathbb{C}^2$, $\ket{\psi}=a\ket{0}+b\ket{1}$ with $|a|^2+|b|^2=1$.
A register of $n$ qubits lives in a $2^n$-dimensional Hilbert space. This exponential size is the
source of both the hopes for QML (large feature spaces) and many of its difficulties (data loading,
vanishing gradients). Mixed states, which describe noise or partial knowledge, are written as density
matrices $\rho$ with $\rho\succeq0$ and $\Tr\rho=1$.

\paragraph{Measurement.}
Measuring $\ket{\psi}$ in the computational basis returns outcome $i$ with probability
$p(i)=|\braket{i|\psi}|^2$ (the Born rule), and the state collapses. What a quantum model outputs in
practice is an \emph{expectation value} of an observable $O$,
\begin{equation}
  \langle O\rangle_\rho=\Tr(O\rho),
\end{equation}
estimated from $M$ repeated measurements (\emph{shots}) with statistical error of order
$1/\sqrt{M}$. Every quantum model therefore pays a sampling cost that classical models do not.

\paragraph{Why measurement matters for learning.}
A measurement extracts at most $n$ classical bits from $n$ qubits and destroys the superposition.
Information that lives in phases and correlations of the state, which is most of it, is lost unless
the measurement is chosen with the question already in mind. This single fact explains why the
\QC/\QQ\ distinction matters (\cref{sec:qq}).

\paragraph{Quantum data.}
``Quantum data'' means data that \emph{is} a quantum state or process: the output of a quantum
simulation, the state stored in a quantum sensor or memory, or a state arriving over a quantum
network. It is different from classical data that happens to describe a quantum system (a table of
measurement outcomes is classical data).

% ===========================================================================
\section{The taxonomy at a glance}
% ===========================================================================

\begin{figure}[h]
\centering
\begin{tikzpicture}[
  cell/.style={draw=#1,line width=1pt,fill=#1!7,minimum width=6.4cm,minimum height=3.6cm,
               align=flush left,text width=6cm,inner sep=6pt,rounded corners=3pt},
  head/.style={font=\bfseries\small,align=center}]
  \node[cell=ccC] (cc) at (0,0) {{\Large\CC}\\[2pt] \textbf{Classical data, classical computer}\\[4pt]
     \small Quantum-inspired algorithms\\ Tensor networks\\ Dequantisation};
  \node[cell=cqC] (cq) at (6.8,0) {{\Large\CQ}\\[2pt] \textbf{Classical data, quantum computer}\\[4pt]
     \small Variational circuits, quantum kernels\\ HHL, quantum PCA\\ \textit{Core of the QML Lab}};
  \node[cell=qcC] (qc) at (0,-3.9) {{\Large\QC}\\[2pt] \textbf{Quantum data, classical computer}\\[4pt]
     \small Phase classification\\ Neural tomography, classical shadows\\ QEC decoders, calibration, control};
  \node[cell=qqC] (qq) at (6.8,-3.9) {{\Large\QQ}\\[2pt] \textbf{Quantum data, quantum computer}\\[4pt]
     \small Quantum convolutional networks\\ Learning with quantum memory\\ Quantum metrology};
  \node[head] at (0,2.2) {Classical processor};
  \node[head] at (6.8,2.2) {Quantum processor};
  \node[head,rotate=90] at (-3.75,0) {Classical data};
  \node[head,rotate=90] at (-3.75,-3.9) {Quantum data};
\end{tikzpicture}
\caption{The four quadrants of quantum machine learning. Rows: what is being learned about.
Columns: what does the computing.}
\label{fig:quadrants}
\end{figure}

The two axes separate two questions that are often blurred: \emph{what are we learning about?}
and \emph{what are we computing with?} \Cref{tab:summary} summarises the quadrants; the following
sections explain each.

\begin{table}[h]
\centering\small
\begin{tabularx}{\linewidth}{@{}l>{\raggedright}X>{\raggedright}X>{\raggedright\arraybackslash}X@{}}
\toprule
 & Typical question & Needs quantum hardware? & Status of advantage claims \\
\midrule
\CC & Can quantum ideas improve classical algorithms? & No & Some proposed quantum speed-ups were matched classically \\
\CQ & Can a quantum computer learn better from ordinary data? & Yes (NISQ or fault-tolerant) & Debated; data loading and trainability are open problems \\
\QC & Can classical ML help us understand or control quantum systems? & Only as the data source & Already useful in practice \\
\QQ & Can a quantum computer learn directly from quantum states? & Yes, plus a quantum data source & Strongest theoretical case; proven separations exist \\
\bottomrule
\end{tabularx}
\caption{Summary of the four quadrants.}
\label{tab:summary}
\end{table}

% ===========================================================================
\section{CC: classical data, classical processing}
\label{sec:cc}
% ===========================================================================

At first sight \CC\ is just ordinary machine learning. It belongs in the scheme because of
\emph{quantum-inspired} methods: classical algorithms that borrow mathematics or ideas from quantum
physics. It also serves a second role: \CC\ is the baseline every quantum claim must beat.

\subsection{Tensor networks}
Tensor networks were developed to represent many-body quantum states compactly \citep{orus2019}. A
state of $N$ particles has $d^N$ amplitudes, but physically relevant states often have limited
entanglement and can be written as a chain of small tensors, a \emph{matrix product state} (MPS),
whose size grows only linearly in $N$.

Stoudenmire and Schwab \citeyearpar{stoudenmire2016} applied the same structure to classification. Each
input feature $x_j\in[0,1]$ is mapped to a two-component vector, and the features are combined with a
tensor product:
\begin{equation}
  \phi(x_j)=\begin{pmatrix}\cos\frac{\pi}{2}x_j\\[2pt] \sin\frac{\pi}{2}x_j\end{pmatrix},\qquad
  \Phi(x)=\phi(x_1)\otimes\phi(x_2)\otimes\cdots\otimes\phi(x_N),\qquad
  f^{\ell}(x)=W^{\ell}\cdot\Phi(x).
\end{equation}
$\Phi(x)$ lives in a $2^N$-dimensional space, like an $N$-qubit product state. The weight tensor
$W^\ell$ for class $\ell$ is stored as an MPS with bond dimension $\chi$, so evaluation and training
cost grows only linearly in $N$ and polynomially in $\chi$. The bond dimension controls how much
``entanglement'', that is how many feature correlations, the model can represent.

\subsection{Dequantisation}
Several early \CQ\ algorithms promised exponential speed-ups for linear-algebra tasks. Kerenidis and
Prakash \citeyearpar{kerenidis2017} proposed a quantum recommendation system whose running time is
polylogarithmic in the size of the preference matrix. Its input model assumes the data can be
prepared as quantum states efficiently, for example from a quantum random access memory
(QRAM) \citep{giovannetti2008qram}.

Tang \citeyearpar{tang2019} showed that a \emph{classical} algorithm given the analogous access, namely the
ability to sample rows and entries with probability proportional to their squared norm
(\emph{$\ell_2$-norm sampling}), can also produce recommendations in time polylogarithmic in the
matrix dimensions, for low-rank matrices. The exponential gap disappears once both sides get the
same input assumptions. Chia et~al. \citeyearpar{chia2022} turned this into a general framework that
dequantises a broad class of low-rank quantum linear-algebra algorithms.

\subsection{Why CC matters for the rest}
\begin{itemize}
  \item Speed-up claims are only meaningful against the \emph{best} classical method with the
        \emph{same} access to the data. Dequantisation is the clearest example.
  \item Every \CQ\ experiment needs a \CC\ baseline with a comparable number of parameters,
        tuned with equal care \citep{bowles2024}.
\end{itemize}

\begin{inshort}
\CC\ is classical computing with quantum mathematics. Tensor networks give compact models;
dequantisation shows that some promised quantum speed-ups came from input assumptions rather than
from quantum mechanics.
\end{inshort}

% ===========================================================================
\section{CQ: classical data, quantum processing}
\label{sec:cq}
% ===========================================================================

This is what most people mean by quantum machine learning and what most PennyLane and Qiskit
tutorials teach: classical data (images, tables, time series) is encoded into qubits and processed
by a quantum circuit. A \CQ\ model has three stages,
\[
  x\ \xrightarrow{\ \text{encode}\ }\ \ket{\psi(x)}\ \xrightarrow{\ \text{process}\ }\
  U(\theta)\ket{\psi(x)}\ \xrightarrow{\ \text{measure}\ }\ \langle O\rangle ,
\]
and usually a classical optimiser that adjusts $\theta$.

\subsection{Data encoding (feature maps)}
\label{sec:encoding}
The \emph{feature map} $x\mapsto\ket{\psi(x)}$ decides what the model can learn, because the
trainable part can only separate what the encoding has already pulled apart \citep{schuld2019feature}.
Common choices:

\begin{description}[style=nextline,leftmargin=1.2em]
  \item[Basis encoding] A bit string $x\in\{0,1\}^n$ becomes the basis state $\ket{x}$. Natural for
        binary data; one qubit per bit.
  \item[Angle encoding] Each feature is a rotation angle,
        $\ket{\psi(x)}=\bigotimes_{j=1}^{n}R_Y(x_j)\ket{0}$. One qubit per feature, cheap to prepare,
        but a product state with no entanglement between features.
  \item[Amplitude encoding] A normalised vector becomes the amplitudes,
        \begin{equation}
          \ket{\psi(x)}=\frac{1}{\lVert x\rVert}\sum_{i=1}^{N}x_i\ket{i},
        \end{equation}
        using only $\lceil\log_2N\rceil$ qubits. Very compact, but preparing a general state needs
        circuits of size $O(N)$, and the norm $\lVert x\rVert$ is lost.
  \item[IQP / ZZ feature maps] Hadamards followed by diagonal gates whose phases depend on
        single features and on products of features \citep{havlicek2019}:
        \begin{equation}
          U_{\Phi}(x)=\exp\!\Big(i\sum_{S}\phi_S(x)\prod_{j\in S}Z_j\Big)H^{\otimes n},\qquad
          \phi_{\{i\}}(x)=x_i,\quad \phi_{\{i,j\}}(x)=(\pi-x_i)(\pi-x_j),
        \end{equation}
        usually applied more than once. These states are believed to be hard to simulate
        classically, which is what makes the corresponding kernel (\cref{sec:kernels}) interesting.
\end{description}

\subsection{The data-loading problem}
\label{sec:loading}
Many theoretical \CQ\ speed-ups assume that a data vector can be turned into a quantum state in time
polylogarithmic in its length. For arbitrary classical data that needs a QRAM \citep{giovannetti2008qram},
which does not exist at scale, or a state-preparation circuit whose cost can be linear in the data
size and cancel the speed-up. Aaronson \citeyearpar{aaronson2015} lists this as the first item of ``fine
print'' to check in any quantum learning claim, together with the condition number of the problem,
the fact that the output is a quantum state rather than a list of numbers, and comparison with the
best classical algorithm.

\subsection{Variational quantum circuits}
A variational quantum circuit (also called a parameterised quantum circuit or quantum neural
network) defines the model
\begin{equation}
  f_\theta(x)=\braket{\psi(x)|U(\theta)^\dagger\,O\,U(\theta)|\psi(x)},
  \label{eq:vqc}
\end{equation}
where $U(\theta)$ is a sequence of parameterised rotations and entangling gates (the
\emph{ansatz}) and $O$ an observable such as $Z$ on the first qubit \citep{mitarai2018,cerezo2021vqa}.
For classification, the sign of $f_\theta(x)$ plus a bias gives the label. Training minimises a loss
such as the mean squared error with a classical optimiser. Only the circuit evaluations run on
quantum hardware, which makes these models suitable for noisy intermediate-scale quantum (NISQ)
devices \citep{preskill2018}.

\subsection{Gradients: the parameter-shift rule}
Back-propagation is impossible on hardware, because intermediate states cannot be stored or copied.
For gates of the form $e^{-i\theta P/2}$ with $P$ a Pauli operator, the exact derivative follows
from two additional circuit runs \citep{mitarai2018,schuld2019gradients}:
\begin{equation}
  \frac{\partial f}{\partial\theta_k}
  =\frac12\Big[f\big(\theta+\tfrac{\pi}{2}e_k\big)-f\big(\theta-\tfrac{\pi}{2}e_k\big)\Big].
\end{equation}
Unlike finite differences this is exact, and it works on hardware. Its cost is two circuit
evaluations per parameter per step, each estimated from shots.

\subsection{Expressivity and data re-uploading}
If a feature enters the circuit through gates $e^{-ixG}$, the model is a truncated Fourier series
in $x$ \citep{schuld2021fourier},
\begin{equation}
  f_\theta(x)=\sum_{\omega\in\Omega}c_\omega(\theta)\,e^{i\omega x},
\end{equation}
where the available frequencies $\Omega$ are set by the eigenvalues of the encoding generators.
Encoding a feature once with a single Pauli rotation gives only frequencies $\{-1,0,1\}$: a single
sine wave. \emph{Data re-uploading} \citep{perezsalinas2020}, encoding $x$ again before each trainable
layer, enlarges the spectrum to $\{-L,\dots,L\}$ for $L$ repetitions. This is why even one or two
qubits can learn curved decision boundaries when the data is re-uploaded.

\subsection{Quantum kernel methods}
\label{sec:kernels}
A kernel method needs only a similarity measure between data points. A quantum kernel uses the
overlap of encoded states \citep{havlicek2019,schuld2019feature}:
\begin{equation}
  \kappa(x,x')=|\braket{\psi(x)|\psi(x')}|^2=\Tr\!\big[\rho(x)\,\rho(x')\big].
\end{equation}
The quantum computer only estimates the Gram matrix $K_{ij}=\kappa(x_i,x_j)$, for example by running
$U_\Phi(x')^\dagger U_\Phi(x)$ and counting how often all qubits return to $\ket{0\cdots0}$. A
classical support vector machine then learns the decision function
$f(x)=\operatorname{sign}\big(\sum_i\alpha_iy_i\,\kappa(x_i,x)+b\big)$. Training is a convex problem,
so there are no barren plateaus in the optimisation itself. Variational models and kernel models are
closely related: the model in \cref{eq:vqc} is a linear model in the feature space of density
matrices \citep{schuld2021book}.

The catch is \emph{exponential concentration} \citep{thanasilp2024}. For expressive feature maps on
many qubits, the off-diagonal kernel values concentrate exponentially around a fixed value, so the
Gram matrix approaches the identity and an exponential number of shots is needed to tell entries
apart. The resulting model memorises the training data and generalises poorly. Huang
et~al. \citeyearpar{huang2021power} add a second caution: given enough data, classical models can often
match quantum kernels, and they give a test (the geometric difference between kernels) for when a
quantum kernel can possibly help.

\subsection{Quantum linear-algebra algorithms}
A separate line of \CQ\ work speeds up the linear algebra underneath many learning algorithms. Such
algorithms usually assume fault-tolerant hardware.
\begin{itemize}
  \item \textbf{HHL} \citep{harrow2009} prepares a state $\ket{x}\propto A^{-1}\ket{b}$ for a sparse,
        well-conditioned $N\times N$ matrix $A$ in time $\tilde{O}(\log(N)\,s^2\kappa^2/\epsilon)$,
        where $s$ is the sparsity and $\kappa$ the condition number. The output is a quantum state;
        reading out all of $x$ would cost $O(N)$ again.
  \item \textbf{Quantum PCA} \citep{lloyd2014qpca} uses copies of a density matrix $\rho$ to apply
        $e^{-i\rho t}$ and phase estimation to reveal the largest eigenvalues and eigenvectors of
        $\rho$ in time logarithmic in its dimension, for low-rank $\rho$.
  \item \textbf{Quantum SVM} \citep{rebentrost2014} solves the least-squares SVM with these
        primitives.
\end{itemize}
Each of these inherits the caveats of \cref{sec:loading}, and several low-rank variants have been
dequantised (\cref{sec:cc}).

\subsection{Trainability and noise}
\begin{description}[style=nextline,leftmargin=1.2em]
  \item[Barren plateaus] For deep, randomly initialised circuits the gradient of the cost averages to
        zero, and its variance vanishes exponentially in the number of qubits
        $n$ \citep{mcclean2018}: $\operatorname{Var}[\partial_kC]\in O(b^{-n})$ for some $b>1$. The
        landscape is flat almost everywhere. Global cost functions (observables acting on all qubits)
        suffer even for shallow circuits, while local costs on shallow circuits keep polynomially
        large gradients \citep{cerezo2021cost}.
  \item[Noise] Gate errors and decoherence shrink expectation values toward their maximally mixed
        values, so deeper circuits carry less signal. On NISQ hardware, circuits must stay shallow.
  \item[Shot noise] Every expectation value, gradient and kernel entry is estimated from finitely many
        shots, so small signals cost many repetitions.
\end{description}

\begin{inshort}
\CQ\ is the most explored quadrant and the one the QML Lab modules teach. Its tools are feature maps,
variational circuits trained with the parameter-shift rule, and quantum kernels. Its open problems
are data loading, barren plateaus, kernel concentration, noise, and showing an advantage over strong
classical baselines.
\end{inshort}

% ===========================================================================
\section{QC: quantum data, classical processing}
\label{sec:qc}
% ===========================================================================

Here classical machine learning is the tool and quantum systems are the subject. The data comes from
quantum experiments or simulations, but the learning runs on ordinary computers. This quadrant is
already productive because it needs no fault-tolerant quantum computer \citep{gebhart2023}.

\subsection{Classifying phases of matter}
Carrasquilla and Melko \citeyearpar{carrasquilla2017} trained ordinary neural networks on spin
configurations sampled from many-body models. The networks learned to distinguish ordered from
disordered phases and located phase transitions without being told the order parameter. The input
is measurement data, that is classical snapshots of quantum states, and the processor is classical.

\subsection{Neural-network quantum states and tomography}
A neural network can represent the wave function itself. Carleo and Troyer \citeyearpar{carleo2017} used a
restricted Boltzmann machine,
\begin{equation}
  \psi_\theta(\sigma)=\sum_{\{h\}}\exp\Big(\sum_ja_j\sigma_j+\sum_ib_ih_i+\sum_{ij}W_{ij}h_i\sigma_j\Big),
\end{equation}
to find ground states of spin models variationally. Torlai et~al. \citeyearpar{torlai2018} turned this around
for \emph{tomography}: training such a network on measurement outcomes reconstructs a many-body state
from far fewer measurements than full tomography, whose cost grows exponentially with system size.

\subsection{Classical shadows}
Huang, Kueng and Preskill \citeyearpar{huang2020shadows} showed how to predict many properties of a state
from few measurements. Each round applies a random unitary $U$ (for example random single-qubit Pauli
bases), measures in the computational basis, and stores a classical ``snapshot'',
\begin{equation}
  \hat\rho=\bigotimes_{j=1}^{n}\big(3\,U_j^\dagger\ket{\hat b_j}\!\bra{\hat b_j}U_j-I\big).
\end{equation}
Averages of snapshots predict $M$ observables to accuracy $\epsilon$ with a number of measurements
that grows only as $\log M$ (for local observables). The measurement is quantum, the learning is
classical: a clean example of \QC, and a common input format for classical ML on quantum data.

\subsection{Decoders for quantum error correction}
An error-correcting code reports a \emph{syndrome}, a pattern of parity checks, from which a decoder
must infer the most likely error. Torlai and Melko \citeyearpar{torlai2017decoder} trained a neural network
to do this for topological codes. Fast, accurate decoders are a bottleneck for fault-tolerant hardware,
and machine-learned decoders are an active area.

\subsection{Calibration, control and characterisation}
Machine learning tunes the hardware itself: optimising control pulses, calibrating gates, and
estimating noise. Bukov et~al. \citeyearpar{bukov2018} used reinforcement learning to find control protocols
for preparing target states. The review by Gebhart et~al. \citeyearpar{gebhart2023} surveys learning for
characterisation, verification and control.

\begin{inshort}
\QC\ uses classical ML to understand and control quantum systems: phase detection, neural-network
states and tomography, classical shadows, decoders and calibration. It works today, at scale, because
the learner is classical.
\end{inshort}

% ===========================================================================
\section{QQ: quantum data, quantum processing}
\label{sec:qq}
% ===========================================================================

In \QQ\ the data is already quantum, and a quantum computer processes it directly without first
measuring it into classical numbers. This avoids the information loss of measurement
(\cref{sec:prelim}) and the data-loading problem of \cref{sec:loading}, because nothing needs to be
loaded.

\subsection{Quantum convolutional neural networks}
Cong, Choi and Lukin \citeyearpar{cong2019} proposed a circuit with the structure of a convolutional network.
\emph{Convolution} layers apply the same quasi-local unitary across the register; \emph{pooling}
layers measure some qubits and use the results to rotate their neighbours, halving the system size;
a final fully connected unitary acts on the few remaining qubits. For $N$ input qubits a QCNN has
only $O(\log N)$ variational parameters. It was shown to recognise symmetry-protected topological
phases directly from the quantum state, and its structure helps avoid barren plateaus.

\subsection{Quantum autoencoders}
Romero, Olson and Aspuru-Guzik \citeyearpar{romero2017} train an encoder $U(\theta)$ that compresses a
family of $n$-qubit states into $n-k$ qubits. The cost asks the $k$ discarded (``trash'') qubits to end
up in a fixed reference state; if they do, the state can be decoded with $U(\theta)^\dagger$ without
loss. Applications include compressing the output of quantum simulations and denoising.

\subsection{Learning with quantum memory}
The strongest results in \QQ\ concern \emph{learning from experiments}. A conventional learner
measures each copy of an unknown state or each use of an unknown process and learns from the
classical outcomes. A quantum-enhanced learner stores copies in quantum memory and processes them
jointly, for example with entangled measurements across two copies. Huang et~al. \citeyearpar{huang2022}
proved that for several tasks, such as predicting properties of physical systems and learning about
unknown quantum processes, conventional learners need exponentially more experiments than
quantum-enhanced ones, and demonstrated the effect on a superconducting processor. This is a provable,
experimentally tested separation, the clearest one in QML so far.

\subsection{Quantum metrology}
Estimating an unknown parameter $\varphi$ (a phase, a field, a frequency) is a learning task on
quantum data. With $N$ independent probes, the error scales as the standard quantum limit
$\Delta\varphi\propto1/\sqrt{N}$; entangled probes and quantum processing can reach the Heisenberg
limit $\Delta\varphi\propto1/N$ \citep{giovannetti2011}. Adaptive, learning-based estimation strategies
connect metrology with machine learning \citep{gebhart2023}.

\subsection{What QQ needs}
\QQ\ needs a quantum data source \emph{connected} to a quantum processor: a quantum sensor, a quantum
simulator or a quantum network link feeding qubits into a processor without measuring them. Such
pipelines exist only in early experiments, which is why \QQ\ is the most promising quadrant in theory
and the least developed in practice.

\begin{inshort}
\QQ\ processes quantum states without measuring them first. It has the clearest theoretical case for
quantum advantage, including proven exponential separations in learning from experiments, but it needs
quantum data sources wired directly to quantum processors.
\end{inshort}

% ===========================================================================
\section{Comparing the quadrants}
% ===========================================================================

\begin{table}[h]
\centering\small
\begin{tabularx}{\linewidth}{@{}l>{\raggedright}X>{\raggedright}X>{\raggedright\arraybackslash}X@{}}
\toprule
 & Key concepts & Main obstacle & Representative work \\
\midrule
\CC & tensor networks, $\ell_2$-norm sampling, dequantisation & none quantum; must beat strong baselines & \citealp{stoudenmire2016,tang2019,chia2022} \\
\CQ & feature maps, variational circuits, parameter shift, kernels, HHL & data loading, barren plateaus, concentration, noise & \citealp{havlicek2019,schuld2019feature,mitarai2018,harrow2009,mcclean2018} \\
\QC & phase classification, neural quantum states, classical shadows, decoders & measurement cost, model expressivity & \citealp{carrasquilla2017,carleo2017,torlai2018,huang2020shadows} \\
\QQ & QCNNs, autoencoders, quantum memory, metrology & connecting quantum data sources to processors & \citealp{cong2019,romero2017,huang2022,giovannetti2011} \\
\bottomrule
\end{tabularx}
\caption{Concepts, obstacles and key references per quadrant.}
\label{tab:compare}
\end{table}

\paragraph{Where does quantum advantage come from?}
\CC\ and \QC\ use classical computers, so they work at scale today. In \CQ\ the case for advantage
rests on feature spaces or linear algebra that are hard to simulate classically, but it is weakened
by data loading, dequantisation and the trainability problems above. In \QQ\ the advantage has a
simple source: measuring destroys information, and a quantum processor can avoid measuring.

\paragraph{Is advantage the right goal?}
Schuld and Killoran \citeyearpar{schuld2022advantage} argue that chasing speed-up claims has distorted the
field, and that QML should also be studied for what it teaches about learning, quantum models and
their inductive biases. Bowles, Ahmed and Schuld \citeyearpar{bowles2024} benchmarked many published quantum
models on standard tasks and found that well-tuned classical models usually matched or beat them, and
that removing entanglement often did not hurt. Careful benchmarking is part of doing QML well.

% ===========================================================================
\section{Where the QML Lab project fits}
% ===========================================================================

The QML Lab is a ten-week study plan (modules~1--8) built on PennyLane \citep{bergholm2018} and
Qiskit \citep{javadiabhari2024}, plus sixteen open-ended research projects managed in a Django portal.
Its base is \CQ: the modules encode classical data into qubits and train circuits and kernels on it.
Every project also carries a classical baseline, which is \CC\ by design.

\begin{table}[h]
\centering\small
\begin{tabularx}{\linewidth}{@{}l>{\raggedright}Xl@{}}
\toprule
Module & Topic & Quadrant \\
\midrule
1--3 & Qubits, gates, entanglement, parameter-shift gradients & foundations \\
4 & Data encoding: angle, amplitude, ZZ feature maps & \CQ \\
5 & Variational classifier with data re-uploading & \CQ \\
6 & Quantum kernels and SVM, with an RBF baseline & \CQ\ (+ \CC\ baseline) \\
7 & Hybrid quantum--classical network in PyTorch & \CQ\ (+ \CC\ baseline) \\
8 & Barren plateaus and noise & \CQ\ trainability \\
\bottomrule
\end{tabularx}
\caption{The learning modules and their quadrant.}
\end{table}

\begin{table}[h]
\centering\small
\begin{tabularx}{\linewidth}{@{}rl>{\raggedright\arraybackslash}X@{}}
\toprule
\# & Project & Quadrant and reason \\
\midrule
1  & Quantum-kernel intrusion detection & \CQ: network-traffic features, quantum kernel \\
2  & Trainable quantum kernels & \CQ: learned feature map \\
3  & Projected vs fidelity kernels & \CQ; borders \CC\ through classical simulability \\
4  & Re-uploading depth study on MNIST & \CQ \\
5  & QAOA for MaxCut with learned angles & \QC-like: a classical model learns to set up a quantum algorithm \citep{farhi2014} \\
6  & Quantum autoencoder & \QQ: compresses quantum states \\
7  & Quantum GAN for a 1-D distribution & \CQ \\
8  & Hybrid modulation classifier & \CQ \\
9  & Quantum transfer learning & \CQ \\
10 & Quantum reinforcement learning on CartPole & \CQ \\
11 & Quantum reservoir computing & \CQ \\
12 & QCNN for phase classification & \QQ: quantum ground states in, quantum processing \\
13 & Learning from classical shadows & \QC: random measurements, classical estimation \\
14 & Noise-aware training with error mitigation & \CQ\ under realistic noise \\
15 & Barren-plateau mitigation & \CQ\ trainability \\
16 & Simulator vs IBM hardware benchmark & \CQ\ on real hardware \\
\bottomrule
\end{tabularx}
\caption{The sixteen research projects by quadrant.}
\end{table}

Each project has its own living document (\path{qml_projects/<slug>/docs/main.tex}) with background,
research questions, method, planned experiments and a results log, using the same Harvard-style
bibliography as this overview.

\paragraph{Reading the mapping.}
Twelve of the sixteen projects are \CQ, matching the tools and simulators available. For a thesis
with a stronger theoretical case for quantum advantage, projects~6 and~12 (\QQ) are the natural
direction; project~13 (\QC) is the most immediately useful for work on real hardware. Every project
follows the same quality bar: a classical baseline of similar size, several random seeds, scaling
with qubit count, and at least one run under a realistic noise model \citep{bowles2024}.

% ===========================================================================
\section{Glossary}
% ===========================================================================

\begin{description}[leftmargin=1.2em,style=nextline,font=\normalfont\bfseries]
  \item[Ansatz] The fixed structure of a parameterised circuit; only its angles are trained.
  \item[Barren plateau] A cost landscape whose gradients vanish exponentially with qubit number.
  \item[Bond dimension] The size of the indices linking tensors in a tensor network; it bounds the
        correlations the network can represent.
  \item[Classical shadow] A classical record of randomised measurements used to predict many
        properties of a quantum state.
  \item[Data re-uploading] Encoding the same input several times between trainable layers to make
        the model more expressive.
  \item[Dequantisation] A classical algorithm that matches a claimed quantum speed-up under the same
        input assumptions.
  \item[Exponential concentration] Kernel values or expectation values collapsing toward a constant
        as qubits are added, so that exponentially many shots are needed to resolve them.
  \item[Feature map] The map $x\mapsto\ket{\psi(x)}$ from data to quantum states.
  \item[NISQ] Noisy intermediate-scale quantum: today's devices with tens to hundreds of noisy qubits
        and no full error correction \citep{preskill2018}.
  \item[Parameter-shift rule] Exact gradients of circuit outputs from two shifted circuit runs.
  \item[QRAM] Quantum random access memory: a hypothetical device that loads classical data into
        superposition efficiently.
  \item[Quantum data] Data that is a quantum state or process, not a classical description of one.
  \item[Shot] One execution and measurement of a circuit; expectation values are averages over shots.
\end{description}

\bibliographystyle{harvardqml}
\bibliography{references}

\end{document}
